% ================================================================
%  18.8 Découverte d'un passé (原书 p.916–926)
%  源文: _work/chunks/chunk15.md
% ================================================================
\section{发现一段过去}

\subsection{(1) 第一口气息——或确认}   % 18.8.1

\notehead{183}{}

我终于来到这场一个多月前开始的回顾与盘点里最切身私己的部分。还剩要快速巡视的,是这场思
考就\textbf{我自己}教给我的东西。

思考让我发现的第一件事,是一段\textbf{过去}——我的数学家过去;此前我从不曾费心在上面停留
过哪怕一瞬。在那光鲜无事、看似平淡的表面之下,我又一次看见了深处的种种展露出来:它们通常被忽略、被抹去(仿佛一把扫帚悄悄扫过),从我们习惯为自己、为周遭描画的那幅舒舒服服的自觉图像里扫了出去。在那些从未检视过(至少在我的数学生涯里如此)的“涂痕”(或扫出门外的尘屑……)之中,有一种阴险、有时咄咄逼
人的东西:我在同某些朋友的关系里的自负。打一开始,这种自负就取了一种数学精英主义的形状;
它秘而不宣,我毫无察觉,因为我的态度看上去太天经地义了。这种精英主义(或如Chevalley与盖杰所
称的“绩优主义”),想必逐年变硬。它结晶为那种“竞技”的姿态——我是在思考“第一口气”接近
尾声时才终于发现的。在一副规规矩矩的门面底下,这种姿态认可的,是一种占有式的、嫉妒的心性,
指向那些被当作我的“禁猎场”的东西,以及那些凭其出色的品质、我乐于放进禁猎场来的人。

\pmark\origpage{1232}所幸,这些非常“老板”的心性,并没有穷尽1948到1970年间,我同数学世界里的朋友、同事
与学生,同数学本身之间关系的内容——差得远呢。尽管如此,它们构成一条阴险的低音线,去年以
前我从未费神去注意它;直到《收获与播种》的第一部分(“第一口气”)。这场渐进的发现,在“竞技
数学”一节(n\textsuperscript{o}~40)里达到顶点。这一节,在我看来标志着思考中一次质变的时刻。当时我立刻
把它感受为一次\textbf{翻越垭口}——它蓦地为我打开一片新的全景……

再退后一年回望,如今我把1948到1970年间、我在数学家中度过的这段漫长的数学生涯的第一
期,看作一场\textbf{交易}:以我本来所属(人人本属)的“长子权”——去完完整整地(倘若这是我
的选择)经历一场独特而独一无二的历险——换来的,是一碗“扁豆汤”:与一个田园诗般、虚构的
“数学共同体”认同(我希望这认同毫无保留,却从未完全做到……),而那个共同体同时又发放
种种舒服的好处\fnote{1027}{(*)}{这就是前面一条脚注(打字稿第1219页注(*))里谈到的那层暧昧。}。借这幅图像,我并不声称已把那段岁月说尽——它太丰厚,套不进一句现成
的公式。但我觉得,它抓住了第一次思考阶段里首次显形的一个重要侧面。这个侧面,又出现在《收
获与播种》这一部分(事后)取得的名字“愚妄与更新”里。

这第一阶段最切身、最深沉的部分,由最后三个“章”\fnote{1028}{(**)}{当然(如我在《收获与播种》引论里所说明的),这些“章”——把由共同主题或特殊亲缘联系起来的相邻各节归拢在一处——是事后设立的,在写完(那将只是)《收获与播种》第一部分之后。在《愚妄与更新》里,我偶尔把它们当作《收获与播种》的“部”来援引(不要把它们同那五个“部”——“愚妄与更新”等等——混为一谈;1984年2月至今的整个思考,归拢为那五个部)。}:第六至第八:“收获”、“孩子自娱”与“孤独的历险”。在“收获”里,我先同我生命中的某些
时刻(不只是我的数学生涯,这一回不是)重新接上头——一些满蓄更新之力的时刻。仿佛受一股
未知之力、一道隐秘而急迫的声音驱使,我寻找着重拾那些\textbf{天真}的心境\dmark{}\origpage{1233},好跨过那道我
隐约还感到横在前面的门槛。当然,当时我无从预料:接下来还等着我的,是发现我对数学本身
的一种占有式的态度。我继续爬坡,不急不缓,仿佛双脚踩着一条只有它们自己才“看得见”的无形
之路。不必对自己说,我也知道它引我去该去的地方;而脚下的雾,一步一步,渐渐散了。

就这样,我抵达了旅途中这道新的门槛,或者说这道\textbf{垭口}:

\begin{quote}
“……一到那个地方,我就有了那种来到瞭望台的人的感觉:从那里,他望得见自己一路走来的景
色——行走之时,他每一步只能看见其中一角。而此刻,有了这种对广度与空间的知觉,这是一
种解放……”
\end{quote}

刚一越过思考的这个敏感点,它便深化为一场对我自身的冥想。第二天,我已感到需要引入“老板”
与“工匠”的形象——别名“孩子”;这形象,两三年来对我已相当熟悉。但我远没有料到,它在接下
来的思考里将证明何等有用——那时,近两个月来我一直自以为马上就要收尾,好马上回到我的数
学笔记,回到《追寻栈》去!

在构成“孩子自娱”这一“章”的四节里,我重新接上了我同数学关系的某些侧面与波折。这些,近
三年前(1981年7月至12月)我已细细探测过,可此后有的是时间,大半又忘了。这一回,我的宗
旨主要是让自己进入状态,去探测我那次不期然而重返长期数学投入的意义,并试着在如今支配我
生命的两种激情之间——表面看互不相容——“给自己”定位:数学,与冥想。

顺便说,这两股激情的“互斥”,如今在我看来已不像两年前那样严酷了。在《追寻栈》里,
数学思考有时让位于、甚至化作一场多少个人化的思考;在那里,我这个有感受、有情感、有(不只
是数学的)好奇、有命运的存在,不再完全缺席。反过来,在\dmark{}\origpage{1234}《收获与播种》这场对我自身的思考里,思考本身又把我同旧日的数学之爱重新接通,并且这里那里,成了数学思考的引信
\fnote{1029}{(*)}{(5月10日)这些“思考的引信”,至今已结了果:我对若干搁置了十五年的主题的理解,得到了更新。}。

这些共存的可能——乃至共生的可能——这两股我心中求知冲动的不同表达,按事物的本性,或许
注定要相当有限。但无论如何,去年的思考(甚至更早,三年前那场长考)已让我看清:这两股
激情绝不互相敌对,甚至本质并无不同。在思考的最后一部分“孤独的历险”里,我竭力贴身把
握:这两股激情究竟在何处不同,它们各自为我打开的“历险”又如何不同。正是在这个追问
中,我发现了一个明摆着、而我一辈子佯装不知的事实:数学是“\textbf{一场集体的历险}”;我个人的
数学历险,只有系于它所属的那场更宏大的集体历险,才获得意义。

说实话,在“孤独的历险”一节里,起初我只是顺带掠过这个事实;当时的宗旨,倒是要用言词把
另一件事勾勒出来——那件事我素来深知,却始终难以全盘接受:冥想,才是一场\textbf{孤独的}\textbf{历
险}。给一件“已知”的事作表述,这番功夫绝没有白费,远非如此!它深化了这份认识,并且顺带让
我发现(至少对我而言)一个明摆着的新事实:把我系于\textbf{另一场}历险——数学的历险,它倒是
集体的——的那条纽带(当时,我、或者说当时的我之中的某人或某物,本想同它保持距离……)。

地基由此打好了:第二天,在“分裂的确认”一节里,我径直进入我困惑的中心。首先是确认:“老板”
的下注,哪怕他想自欺(那本更合他的天性……),也只能是那场集体历险——唯一能给他“回报”
的历险。“孩子天性孤独”;\dmark{}\origpage{1235}唯有孩子,才会被一场世上无人想要的历险、一种切实可触且常常明摆着
的知识所吸引,而这知识,他无法同任何人分享。而此刻,在我这桩“事业”里,“小
家伙的偏好”,偏偏不合“老板”的意,正落在这一边。

这个确认,引出对我之内一场\textbf{分裂}的确认:\textbf{老板与孩子的分裂}。第一次,我在如此专注、
如此严格的心境里,作出这样一个确认。这不是一道\textbf{判词}——依照这样那样的“看法”或哲
学或随便什么拟出、自命多少具有普遍效力的判词。它只是一个简单的\textbf{确认},出自对一个极特
殊的个案的细心检视:我这卑不足道的人,在我发展的某个阶段。也许有一天,这道分裂会消失,
而老板照旧去做必要的事,任工匠孩子随意做工。这不是我今天操心的事,也不必是。一天的难处一
天当就够了……

(4月5日)确实,这道分裂,差不多九年前,已在一个梦里,借一场以震撼人心的力量搬演出来的
寓言,向我揭示过。那是在发现冥想——这种长期被我无视、随时待命的力量——的两天之后;正
是掘到这个梦的意义的深处,我才重新找到我之内那个不分裂的东西,\textbf{另一个}我——沉默而
不可见了那么久的,“一位至亲之人,被人当作死了一辈子……”(见《愚妄与更新》)。当时出现的
新东西、要紧的东西,并\textbf{不是}那道我熟得不能再熟的分裂,也不是梦以那样大的力量向我揭示
的、这道分裂的本性——它化身成两个亲爱而熟悉的存在,谁也没有名字,而他们本是\textbf{同一
个}——而是那场\textbf{重逢}:在四个钟头高强度冥想(如同分娩的剧烈阵痛)之后来到的重逢。

我当时就知道,此后几天、几星期里也一直知道:重逢并不等于分裂的终结。但托重逢的福,我用
新的眼睛看这道分裂了——它固然要紧,但在另一个更根本的真实面前,终究是“次要的”:那是我
重新找到的我之内那个\textbf{不分裂}的、不可摧毁的统一,我后来认出它就是“孩子”。这一双重的
\dmark{}\origpage{1236}认识,当时鲜活而敏锐地在场。随后几年里,它钝化了——那份对这道“次要”却真真切
切的分裂的认识,渐渐被收藏起来。而当“老板”由着自己一头扎进冥想(那匹出名的“三脚马”……),
他很想暗示(却始终没有胆量、或没有笨到明说……):有了冥想和这一切,分裂从此已成过往,几乎
可以说没有分裂了——顶多这里一小道涂痕、那里一道,好吧,咱们不否认,但总归几乎跟没有一样;
只消看看小家伙工匠多么兴高采烈地埋头干活,看看一团和气的老板怎样踮着脚尖走路,生怕惊扰
了他——真可谓琴瑟和鸣!我不免疑惑:去年那份思考——翻越垭口(“竞技数学”)之前的那份,
尤其当我出人意料地回顾“我的激情”时(同名的节,n\textsuperscript{o}~35)——是不是恰恰还带着这样
的调子:光照逼得人对玫瑰色起一丝疑心……

不管怎样,这道“分裂的确认”当时恰到好处地让我重新接上了一个我多年渐渐看不见的真实。同
时,它让我以新的照明、新的眼睛,重新看见八年前我曾知觉得清清楚楚的那道分裂。我可以毫无保
留、毫无怀疑地这样说,因为我记得清清楚楚:作那个“确认”的时刻,同重逢的插曲、同那插曲就
某道分裂及其本性教给我的东西,没有任何联想!这一联想,直到稍后——我重新拾起前一天的笔
记时——才姗姗而来。这再清楚不过地表明:那次插曲中显现的认识的“次要”(而且不受欢迎!)
的内容,被收藏得多么干净。而且当年做来想必格外容易:重逢那道关键转折之后,没有人就这份
内容作过思考;而最能表达它的那幅形象——“老板”与“工匠孩子”——要再过好几年才出现。

如今我似乎看清:这道重新作出的“分裂的确认”,才是《收获与播种》第一部分里我对自己学
到的最要紧的东西。这个确认,不过几行字,是思考这一部分里最短的几节之一。\dmark{}\origpage{1237}人们也许
会想:既然只为了走到这一步,或许犯不着再用一百五十页,去追索自负在我数学生涯中的种种曲
折显形。就流行的“常识”而论,这千真万确。可同样真确的是:这种用修枝刀裁出来的“常识”,根
本无力把握一场发现工作的幽深曲折——无论是认识自己,还是更粗朴的数学发现的工作\fnote{1030}{(*)}{科学发现的工作,在我看来之所以比认识自己的工作“更粗朴”,(我想)有两个缘故。其一,它几乎只调用我们的智性官能,也就是我们存在中极小的一部分。(而且,科学工作倾向于让这部分官能增生肥大,以其余官能、以人格的整体平衡为代价,及至把人变成一种电子计算机怪物……)其二,科学工作所调动的、对抗揭示真实的内在抵抗,通常与对抗自我认识的抵抗,不成比例。这也是为什么,“科学历险”很少——今日几乎从不——是一场“求真的历险”:一场调用我们谦卑的能力、调用我们承担不期而至之真实的勇气的历险——首先对我们自己,然后对外部世界。}。我有一个切身的信念:在这场《收获与播种》的长考
里,每件事都来在其所、来在其时,由先来的一切所准备、所酝酿。

\subsection{(2) 第二口气息——或调查}   % 18.8.2

\notehead{184}{}

(4月6日)随着这道关于分裂的短短确认——去年3月末(一年多前),我起初以为自己已经写完
了《收获与播种》。我远没有料到,它还要再来五倍之多(按页数算)!随后几天,我做些杂事,思绪
开始回到数学主题。可是,思考中还有一个悬着的“小尾巴”,继续在我脑子里打转。越过一层看似
纯然细节的困惑,我隐约感到:对于“老板”转向一项长期数学投入的背后那股力量,我还没有真正
走完一遍。或者说,即便我把它的主要机关都翻了出来,我的理解仍旧苍白而游移——只因为在这
件事上“摆”得还不够久,没能让它渗得更深。这个“最后的尾巴”,将成为我回过头去看那团模糊印
象的门径。思考的重新接续,发生在当时(并且此后三个星期里)被当作《收获与播种》终篇的那一
节\dmark{}\origpage{1238},它旋即得到名字“一段过去的重量”。这个名字,说尽了那场意外的发现:我的数学家过去的
这份\textbf{重量},以及把我继续系于那场集体历险的纽带之力。而且,那天我所瞥见的,还只是
冰山露出的那一点小小峰尖;那沉在水下的庞大部分,要在随后的几个月、整整一年里,才渐渐
显形……

为思考的第一口气收尾的这一节,同时像第二口气的引信与召唤。这份“过去的重量”,显然扎根于
我对一份著作的眷恋——与其说眷恋那部完成了的、做到终点的著作,不如说眷恋一些执念观念与
一些洞见;它们的丰饶与力量,我感觉得到、我“熟知”入骨,而我多少含混地、多年以来就意识
到:它们却在贫瘠而干旱的土地上萎顿着,那土地暗地里、阴险地怀着敌意……于是,“一段过去的
重量”这场思考——它唤回我的好记忆:那部著作,以及我同它的纽带——成为一道长注的机缘;自
“离去”以来,我第一次就这部著作、就它遭到的命运表明态度。十年来、十五年来含混感到的东西,
终于成形,化作言词——言词有时迟迟不肯来到;而一旦白纸黑字写下,它们清清楚楚地向我道出一
个讯息,直到那时我一直在回避去知悉的讯息。后来,鉴于这条一气写成的长注的篇幅,我把它分成
两道,取名“我的孤儿们”与“拒受遗产——或矛盾的代价”(n\textsuperscript{o}s~46、47)。

\enlargethispage{2pt}
可以认为,这双注,就是\emph{葬礼}思考的发令枪\fnote{1031}{(*)}{可惜,这一情节在“葬礼(I)(或中国皇帝的新衣)”的目录里看不出:那道双注构成送葬行列之二(“孤儿们”),而非行列之一(“身后的学生”)。缘由在于“愚妄与更新”最后一节“一段过去的重量”(n\textsuperscript{o}~50)——那几道注本应是对它的评说——之中,诸“注”(n\textsuperscript{o}s~44至47)引用的先后次序。}。三星期后,4月19日,思考继续推进:那本“难忘文集”LN~900激起的
震动,为 Motif 的掘墓复生、由“新爸爸”Deligne主持,作了证。思考的这“第二口气”,一直高密度
地持续到5月底、6月中——当时我再一次自以为马上要画上句号,货真价实的句号!——以一场病\dmark{}\origpage{1239}
中插曲告终\fnote{1032}{(*)}{关于这场病中插曲,见“事件——或身体与精神”与“陷阱——或便利与衰竭”两条注(n\textsuperscript{o}s~98、99)。}。

这第二口气,严格说不是一场对我自己、对我过去的思考,而毋宁是一场对刚发现的\emph{葬礼}的
“调查”,同时是一场努力:把那些明摆着却(大概鉴于我那拔不掉的天真!)令人瞠目、令人难以置
信的事实,边走边尽量“消化”。不过,它毕竟就我自己教给了我一些东西:主要是让我切身掂出,我
对我的过去、对我的著作的眷恋之力。我被触到了痛处:眼看着著作被一块块撕走——有的进了垃
圾堆,有的供人取笑,有的被肆无忌惮地据为己有,当作谁都可以拿的小玩意……

那时我知道,我并没有像我一度相信的那样(在翻过某个“垭口”、眼前展开广阔全景之后
的那阵欢腾里\fnote{1033}{(**)}{这阵欢腾,表达于“旋转木马收摊了!”一节(n\textsuperscript{o}~41);五六星期后,在注“一脚踏进旋转木马”(n\textsuperscript{o}~72)里被压低了调门。})“离开旋转木马”。换句话说,我得以掂量这份过去的全部\textbf{重量},以及把我系
于其上的那些唯我机制的全体之力。这是个大意外!

然而,在思考的第二阶段,我关于自己还有另一项发现,补足第一阶段的所获。第一阶段,我翻出的
主要是我身上自负态度的某个“反面”:以\textbf{排斥}的姿态,对待某些同事甚至朋友——因为这样
那样的缘故,我不把他们算进我自觉身属(当然,是默而不宣的!)的“精英”的世界。同一枚奖牌
的\textbf{正面},则是我在同更年轻的数学家(尤其是我的学生)的关系里的一种\textbf{纵容}与暧昧:我
可以说把他们“吸纳”进了“我的世界”;或者因为他们手段出色,或者仅仅因为我收他们做了学
生,他们从此便被我知觉为处于我的“庇护”之下。我是在5月10日的注“高升”(n\textsuperscript{o}~63$'$)里开始
点到这种态度的,继而是5月27日的注“格格不入者”(n\textsuperscript{o}~67$'$),两道都写我同我年轻而\dmark{}\origpage{1240}光辉的
朋友皮埃尔的关系。这场思考在6月1日的注“暧昧”(n\textsuperscript{o}~63$''$)里深化,转向我同
我的学生们的一般关系。正是在那里,我终于察觉某种暧昧:它从未被我察觉、检视过,一直跟着我
到近几年。而且,就在最近,我又在稍有不同的语境里同这种暧昧打了个照面:在子注“驱逐(2)”
(n\textsuperscript{o}~169$^{1}$)(其第二部分,日期为3月16日)里。我不得不在那里确认:把我本人从 SGA 讨论班
中排挤出去——那个讨论班,是我十年生命投入的总和\fnote{1034}{(*)}{关于这场排挤最新的一桩,见注“殡仪馆——‘为科学服务’”(n\textsuperscript{o}~175)。}——这一排挤,主要由我旧日学生中最亲近的某几
位操办,不过是我在他们面前乐于经营的那种暧昧态度的自然收成:关于他们在那部规模恢弘的
SGA 里、与我在其中的各自位置——他们中的某一位,曾在其中投入一两年光阴。

\subsection{(3) 第三口气息——或发现暴力}   % 18.8.3

\notehead{185}{}

还剩要巡视的,是思考的“第三口气”就我自己教给我的东西。第三口气始于去年9月22日(病中
插曲之后),眼下正接近尾声\fnote{1035}{(**)}{《收获与播种》的第五部分,我在这里暂且撇开;它最初本是葬礼仪式内部——甚至就在“阴阳之钥”内部——的一段“枝蔓”。写这篇回顾时,这一部分尚未完成,故未收入。(6月22日)随后几星期里事情显明:“阴阳之钥”之后的“四种运算”部分(即葬礼(3)),构成《收获与播种》的“第四口气”,同样不在这篇最终回顾之列。}。这里首先指的,是在“阴阳之钥”里继续的那场思考;在我看来,
它是《收获与播种》最切身、最深沉的部分。没有任何刻意的安排,我的这个人、我与世界的关系,
在那里最常处于注意的中心。当注意力似乎偶尔岔开,转向一些表面上更一般的主题,或在我朋友皮埃尔其人身上流连,它却总是回到中心、回到那位演员兼观察者、那个感受着、知觉着、追问着、探
测着的人身上,仿佛被一股无形之力吸回。所以,先于一切、尽管并不以之为目标,它是一场\textbf{对我
的生命、对我自己的冥想}——从一个意外的入口进入:那场\emph{葬礼}。

这也是思考中在我看来最丰饶的部分,是我从中获益最多的部分。许多“已知”的东西彼此各安其
位\dmark{}\origpage{1241};一些仅仅瞥见或预感到的东西,一些“知道”却被忽视、被淹没在日常杂沓的昏暗光晕里的东西,
开始从暗处浮出,既显出它们的重量,也显出它们的轮廓。像是一次新的敞开,一次向未知再次出发
的邀请——而在此刻,那场有名的“长期数学投入”眼看要让认识自己的工作一放就是好几年……

在这里,不适合原原本本地重述这场长考的各个阶段,也不适合为它教给我的一切开一张“清
单”。我宁可几句话说说:对认识我自己而言——作为一场仍在日复一日、月复一月、年复一年继续
的成熟的材料——什么在我看来最要紧。

这场思考开始时,我本意是开一个“括号”(顶多一两道注的工夫……),好让读者(顺便也让我自
己)“入水”,浸进一种阴阳(或“阴柔-阳刚”)的事物辩证洞见。开这个括号的缘由,是需要用阴
阳的直觉,去定位一个触目的印象:它来自对某一篇“悼词”的检视\fnote{1036}{(*)}{关于这篇“悼词”(其恭维之语,调配得法,施用有道……),见同名的两条注(n\textsuperscript{o}s~104、105),以及部分概括了它的注“宝石们”(n\textsuperscript{o}s~170(iii))。}:那是一种刻意的“角色对调”,发
生在一场原本的阴阳关系里。括号于10月2日打开。直到11月10日——在关于阴阳的把戏在我
的生命里、在一般的存在里、最后在数学发现的把戏里的整整一百页密不透风的思考之后——时
机才算成熟,可以至少把这个六个月前冒出的联想\textbf{说出来}\fnote{1037}{(**)}{在注“对调(3)——或阴葬阳”(n\textsuperscript{o}~137)里。};至于在知根知底的状态下去探测它,还要再过十四
天\fnote{1038}{(***)}{在注“天鹅绒的脚掌——或微笑”(n\textsuperscript{o}~137)的开头。}。(而著名的阴阳括号,要到将近两个月后、1月14日,才终于合拢——甚至过了一阵,我才察觉它早已合上了……)

\pmark\origpage{1242}很快,无须寻找也无须计划,“冲突”——在人的生命里、在人身上的冲突——占据了注意的
中心。\emph{葬礼}的发现骤然而强力地调动起来的唯我能量,在那里充当了一支援兵,让我再一次、
就在当下,直面多年以来一直向我发问的那个“冲突之谜”\fnote{1039}{(*)}{这一“发问”,主要开始被察觉,是在我1979年8月至1980年3月间持续的那场对我父母的漫长冥想之中。}。早在此前那些年里,这个谜已渐渐浮到我
想探测、想理解之事的前排,能走多远走多远——却从来没有“跨出那一步”、把自己整个儿投进去
……

渐渐地,思考之中显形出来的,是我生命里仿佛那枚“硬核”的东西,这个谜那可怕的震中,仿佛
“恶之谜”的心脏本身:一种可以称作“无缘无故”或“无因”的暴力——一种似乎单以伤人、害人、
毁人为快的暴力,一种从不报出姓名的暴力,常常闷声不响,顶着一副天真无辜、和蔼可亲的面孔,
因而在触及与摧残上愈发有效——“天鹅绒里的利爪”,纤细、敏捷、毫不留情……在“天鹅绒里的
利爪”那组注(n\textsuperscript{o}s~137--140)的思考中,注意最终落到的,正是这种暴力;直到《阴阳之钥》
终卷,它也一直处于注意的中心。它构成那把最高音:最后一道注里,那条“业”的“没有尽头的锁链”
——父母传给子女,子女传给孙辈,一代又一代,自太初以来——被唤起。

这是我一生中第一次,直面这个“无恨亦无情”的暴力之谜——一种深植于人类生活的暴力,自我
少年时代起,就在我生命里打下了洗不掉的印记。这也是我第一次,在我自己的存在里,确认这道印
记。这同时也是对一个简单事实的确认:这种暴力的\textbf{存在},它那可怕的无所不在——在我的生
命里,也在每一个人的生命里\fnote{1040}{(**)}{这一确认,构成注“无恨亦无容赦”(n\textsuperscript{o}~157)那场思考的高峰时刻。}。而这一个简简单单的确认,本身就在胚胎里含着对这个可怕事实的一
种\textbf{接受}。或许正是在这个确认里,藏着整部《收获与播种》的思考中我所\dmark{}\origpage{1243}学到的最重要的东西
(或者说,至少\textbf{开始}学到的东西)。

这并不是一个终局,不是一场思考的顶点。毋宁说,这仍是第一步,把我带过一道通向未知的门
槛。对我的行路与成熟而言,这个谦卑的步子,在我看来,分量比我所瞥见的那些“回答”的胚胎更
大——在随后几天里,我瞥见了“无因的暴力”之“因”这一问题的答案的胚胎\fnote{1041}{(*)}{见同名之注(n\textsuperscript{o}~159)。}。这个问题本身,只有
在充分看见并承认为其发问的那个事实的存在与分量之后,才获得它的全部意义——那分量,远比
一个“心理机制”的问题沉重。

有人会说,我这是跑题了:对一个(或我自以为)关乎人人的普遍心理事实的确认,属于客观知识的
领地,那是各门科学的专属(心理学、精神病学、社会学,或随便什么),不属于那出了名的“自我认
识”的领地(一种被感觉为含混而不可捉摸、倘若不算全然想入非非的领地)。可我看得很清楚——不
是含混不可捉摸地看,而是像看一个熟悉的、明摆着的数学事实一样清楚……:脱开认识自己,这样
的确认就失去了它活的意义——失去那使它不至于沦为一场哲学-心理学的文体练习、沦为一条
“论纲”的展开(当然十分有意思,如此等等……)的东西。这个确认,本身是一场\textbf{发现},一场贴
身的发现:世上没有任何人能替我作这个发现,我也无法替世上任何别人作它。这场发现,是那次发
现自己的旅程中,最近的一步——或者几乎是最新的。它把我相对于一件要紧而可怕的事定位:这件
事打过我烙印,而我直到那时一直执意忽视它,仿佛我是由于某种特殊的倒霉(也许系于我这卑不
足道的人的某些特点),才一生暴露于它;仿佛别人暴露于它、或施它于人——只要我肯费心睁眼看
看周围——也无非同样的倒霉。

而且这肯定不是偶然:对暴力的思考刚一起头,思考自身的内在逻辑,就引着我——也是平生第一
次——回头去看我记得的那几桩个案:我自己也曾\dmark{}\origpage{1244}\textbf{施}这种“超出理解的暴力”于人,当然,不
曾多想\fnote{1042}{(*)}{见《阴阳之钥》里“天鹅绒里的利爪”部分之后的那道注“正义者的暴力”(n\textsuperscript{o}~141)。}。这次回望的意义,不在于给我机会自责(还是在众目睽睽之下!)——这种事,我反倒完
全没有去做。而在于它为我打开了一扇门,通向对暴力更深的理解——一扇门,往后的日子里,何时
跨过它,只由我定。

\subsection{(4) 忠诚——或女性气质的数学}   % 18.8.4

\notehead{186}{}

这就是,以发现自己的旅程的眼光看,在我看来最要紧的东西。阴阳思考的最后阶段——以暴力为
中心——贯穿最后四部分:“天鹅绒里的利爪”、“暴力——或游戏与刺棒”、“另一个自己”、“冲突与
发现——或恶之谜”,从12月7日到1月14日(它们略多于《阴阳之钥》的三分之一)。

退后看去,我似乎觉得:《阴阳之钥》前八部分的主要角色,是最终把我带到这场关键的思考。我在
这准备性的部分里展开的许多东西,是多年以来我熟悉的事;却仍须“重提”,好让一位初来的读者
能够跟上,好给思考一种内在的融贯——否则它就有缺失之虞,或只有我一个人看得见。不时地,文
体显出这些内在心境的痕迹:写的人急着赶紧把重提打发完,好终于抵达“正题”——虽然,那些所谓
的重提,常常分量重得多,值得我在上面多停留片刻,比我那么急着要去的那个“正题”更值得(而
无论急不急,我到那“正题”,还是晚了一个多月……)。这些心境,我看在三节相连的部分里最为
可触:“夫妻”、“我们的母亲死神”、“拒与受”。即使在那里,真的,重新接上那些本该“已知”的事
情时,我仍旧不由自主同时重新认出了它们,而且有时是在新的光线下——哪怕是为了一些初看如
此非个人化的事情,比如清点那些由直接的亲缘联起来的阴阳“夫妻对”(或“锁孔”)——那些通往
世界的门。

但正是在接下来的三部分里(它们也正好先于以暴力为中心的最后四部分),我\dmark{}\origpage{1245}重新驶入此前未曾
勘探的岸域:“阴阳数学”、“阴阳的对调”、“主人们与仆人们”。

正是在这三部分的第一部分里,发生了那场“大意外”:它将为\emph{葬礼}的意义——或至少某种意
义——投下新的光。事情是这样的:在我对待数学的进路里,更一般地,在我走向发现世界的自发举
步里,我的存在的基调是\textbf{阴,“阴柔的”}。换个说法:尽管我的自我的条件化结构、我这桩事业
的“老板”是阳(不说“彻头彻尾的男子汉”吧),我本来的天性、我之中的“孩子”(他同时是那个
把孩子在玩耍中发现的东西做出形状的工匠……)却以“阴柔”为主。当然,单是这个特点,并不足
以把我个人对待数学的“风格”同任何别人区分开来。我甚至觉得,即使在数学家当中,这个基调(或
“主导”)本是阴的,也并不那么罕见。而我的情形里反常的(在我看来)倒是:在我的发现步履里、
尤其在我的数学工作里,我一生都对这本来天性完全忠诚,从不曾起过一丁点儿要修改它、纠正它的
念头——无论是出于内心检查员的意愿(反正他从来什么也没看见:在一门“男人之间”的事情里,
谁会怀疑一种“阴柔”的感受力与创造性的进路呢!),还是出于迁就外部世界、尤其是科学世界里流
行的好品味的法度。我毫不怀疑:主要正是靠着对我自己天性的这份忠诚——至少在我生命中这个
有限的领域里\fnote{1043}{(*)}{如我在《收获与播种》各处有机会一再说起的,自我八岁起(直到1976年,我四十八岁)支配我生命的两股最强的唯我之力之一,是对我身上“阴柔”特质的压抑,以利于那些被感受为“阳刚”的特质。只是到了“阴阳之钥”的思考里,我才意识到:这种压抑,并没有施行在我的数学工作上(后来,也没有施行在冥想、认识自己的工作上)。我存在中本来“阴柔”的主导,得以在一项通常被感受(也难怪)为最“阳刚”的活动里,尽情尽兴!(关于这一点,见注“最‘男子汉’的艺术”,n\textsuperscript{o}~119。)},我的数学创造力才
得以充分地、无阻碍地铺展,像一棵茁壮的树,稳稳扎根在大地里,依着昼夜、风与季节的节律自由
舒展。事情是这样成全的,尽管我的“天分”颇为平平,尽管开端根本算不上\dmark{}\origpage{1246}\textbf{吉兆}\fnote{1044}{(*)}{我说“天分平平”,绝非故作谦虚。这件事我一遍遍有机会确认:无论是同那些光辉的数学家打交道——他们把握要领、接受并消化新思想,比我敏锐得多——还是同这样一些没有名姓、没有像样数学底子的学生打交道,他们的好奇心与数学发明力一时被唤起。
我的“开端”(至少是我同数学家世界接触的开端,1948年),我在“受欢迎的异乡人”一节(n\textsuperscript{o}~9)里谈了一点。不过,还要早三年,1945年,我的“数学生涯”就开始了:我的精力的大半,投给了数学研究工作。直到1949、1950年前后,作为在法国的一个外国人,我以数学谋生的前景,还成问题得很。倘若这样的机会没有出现,我打算去学木匠——那是一门合我口味的营生。}。

当我在注“上涨的海……”(n\textsuperscript{o}~122)里,对我的数学进路作出这项出人意料的确认时
\fnote{1045}{(**)}{另见后文注“箭与浪”(n\textsuperscript{o}~130)。},它来得有点像一桩意外的奇事,有点“在边上”——因为在我的生活里,同他人的
关系,全都打着我阳的、超级阳的抉择的印记。只是在这场以冲突的动势为中心的思考的继续里,
借着对\emph{葬礼}的一次回望,我才意识到:我的数学家同侪同我的这个人、尤其同我的著作的关
系,被打上了这个不寻常的特点何等深的烙印——在他们心里唤起矜持的反射(若不是排斥),面对
一种被暗暗感受为“不合适”(不说“失体”)的进路风格。这类反应,早在我初入数学世界时就很有
些;只是在那宽厚的年月,它们被当时盛行的敬重他人的氛围调和着——至少在我落脚的那些数学
圈子里。后来,鉴于“Grothendieck结果之力”(借Borel致Mebkhout信中的话,那封信里提到了这些
“矜持”),它们也只好被压着。而自我告别数学舞台以来,它们成了常规,有时还在某副得体的低声
调门(那种低声调门,至今仍是规矩)背后舒舒服服地铺展:昔日的敬意早已风化、消失多时,而当事
人(据推定已死且已葬)也不在场,不能接话了……\emph{葬礼}的这个出人意料的面向——以我这卑
不足道的人为象征,埋葬“数学中的阴柔”——在“天赐的机缘——或封神”与“拒认——或召回”两
道注(n\textsuperscript{o}s~151、152)里得到探测;那是12月23日与24日,正在关于暴力的冥想的正当中。

\pmark\origpage{1247}还剩我这个人最后的一个面向想谈一谈——它是在写《阴阳之钥》时、在刚引的最后那一部
分“主人们与仆人们”里(那一部分紧挨着由“天鹅绒里的利爪”开启的思考转折之前)显形的。这
就是“服务的冲动”,以及它在我对数学投入的抉择里扮演的头等角色,以及它作为一股活的力,在
那些浩大而没有尽头的奠基任务里运转——我之后,还没有别人找到勇气(或谦卑……)去接手并把
它们继续下去。这个以非凡之力在我身上在场的面向,雄辩地见证了我本来天性中“阴柔”的主导:
它保存了下来(乃至避难……)在数学活动里(谁也不会想到去那里找它……)。

此刻我心中闪过一个念头:说不定,正是这股冲动,也以它非唯我的那一份,参与了那次“转向”——
转向高强度的数学活动,而把冥想的功课撂到了后台,一放就是不知多久。冥想,按其本性,是一项
孤独的功课;一项(在我看来)除非自欺、便无法纳入“为众人服务”、为某个“渴望知道的存在的理
想共同体”服务的投资眼光的功课。这么看来,似乎有一股深层的冲动,不同于唯我的求证之心、求
许之心,一股表达着个人与其所属的种属之间的深层纽带的冲动——它注定要在一场旷日持久的冥
想功课里受挫,按我对冥想的理解。而这也许是一条附加的缘由——除了那些来自唯我结构(“老
板”的心性)的缘由(它们本身已经足够强大)——使得这样的功课显得如此稀罕,以至于我拿不
准,可曾在别人身上见过它的痕迹。
